% Data preprocessing flow diagram
\begin{figure}[!t]
\centering
\begin{tikzpicture}[
    node distance=0.45cm,
    stepbox/.style={rectangle, rounded corners, draw=black, thick, minimum width=6.6cm, minimum height=0.65cm, align=center, font=\scriptsize, fill=green!8},
    arrow/.style={-{Latex[length=2mm]}, thick}
]

\node[stepbox] (s1) {1. Column detection (date / target / group / exogenous)};
\node[stepbox, below=of s1] (s2) {2. Date parsing cascade \& target numeric coercion};
\node[stepbox, below=of s2] (s3) {3. Index set, sorted; duplicate timestamps summed};
\node[stepbox, below=of s3] (s4) {4. Frequency inferred; calendar gaps reindexed};
\node[stepbox, below=of s4] (s5) {5. Missing values: interpolate $\to$ ffill $\to$ bfill};
\node[stepbox, below=of s5] (s6) {6. IQR outlier detection and capping};
\node[stepbox, below=of s6] (s7) {7. Exogenous features aligned (price: mean/ffill; promo: max/zero)};
\node[stepbox, below=of s7, fill=blue!10] (s8) {Clean target series + aligned exogenous frame};

\foreach \a/\b in {s1/s2, s2/s3, s3/s4, s4/s5, s5/s6, s6/s7, s7/s8}
    \draw[arrow] (\a) -- (\b);

\end{tikzpicture}
\caption{Automated data preprocessing pipeline applied to every uploaded dataset prior to model training, exactly as executed for the evaluation dataset in Section~\ref{sec:dataset} (Table~\ref{tab:preprocessing-steps}).}
\label{fig:data-flow}
\end{figure}
